\BeleskeID{09_3-s002}
% element: 09_3-s002:e03
% element: 09_3-s002:e04
% slide-id: 09_3-s002
Za formiranje kanala potrebno je dovesti naelektrisanje na gejt. Električno polje gejta privlači iz osnove nosioce koji obrazuju kanal. Punjenje i pražnjenje ove i drugih kapacitivnosti traje konačno vreme, pa one dominantno utiču na brzinu rada.

Invertor pri promeni izlaza puni ili prazni sopstvene izlazne kapacitivnosti i ulazne kapacitivnosti narednog stepena. Istovremeno, njegov ulaz predstavlja kapacitivno opterećenje prethodnog invertora. Kapacitivnosti vodova postoje, ali ih u ovom razmatranju uzimamo kao relativno konstantne.

Promenom dimenzija podešavali smo statičku karakteristiku prenosa. Ista promena menja interne i ulazne kapacitivnosti, pa statičko dimenzionisanje utiče i na prelazni režim.
